% Vietnamese translation for OI Public Library
% Source-PDF: IOI中国国家候选队论文/2000/肖洲论文.pdf
\documentclass[11pt]{article}
\usepackage{oipl-vn}

\title{Ứng dụng cấu trúc dữ liệu trong lập trình}
\author{Xiao Zhou}
\date{}

\begin{document}
\maketitle

\origtitle{Applications of data structures in programming}
\sourcepath{IOI China candidate team papers / 2000 / Xiao Zhou.pdf}

\renewcommand{\contentsname}{Mục lục}
\tableofcontents

\paragraph{Từ khóa.}
Cấu trúc lôgic; cấu trúc lưu trữ; tối ưu hóa thuật toán.

\section*{Tóm tắt}

Cấu trúc dữ liệu là nền tảng của lập trình, nên ảnh hưởng của nó đối với hiệu
suất thuật toán hiển nhiên không thể xem nhẹ. Bài viết này bàn về cách chọn
cấu trúc dữ liệu hợp lý để tối ưu hóa thuật toán, đồng thời thảo luận một số
nguyên tắc và phương pháp lựa chọn cấu trúc dữ liệu.

Trước hết, bài viết phân tích tầm quan trọng của cấu trúc lôgic của dữ liệu
và nêu hai nguyên tắc cơ bản khi lựa chọn cấu trúc lôgic. Tiếp đó, bài viết
so sánh ưu nhược điểm của hai cấu trúc lưu trữ tuần tự và liên kết, rồi bàn
về cách chọn cấu trúc lưu trữ. Cuối cùng, từ một góc nhìn khác của việc chọn
cấu trúc dữ liệu, bài viết tiếp tục thảo luận cách kết hợp nhiều cấu trúc dữ
liệu với nhau. Trong khi bàn phương pháp, bài viết cũng chọn một số đề thi
tin học có tính đại diện để phân tích làm ví dụ.

\section{Dẫn nhập}

``Cấu trúc dữ liệu \(+\) thuật toán \(=\) chương trình''. Câu nói ấy cho thấy
bản chất của lập trình là chọn một cấu trúc dữ liệu thích hợp cho bài toán đã
xác định, rồi thiết kế một thuật toán tốt. Vì vậy, cấu trúc dữ liệu giữ vị trí
rất quan trọng trong lập trình.

Cấu trúc dữ liệu là tập hợp các phần tử dữ liệu mà giữa chúng tồn tại một
hoặc nhiều quan hệ đặc thù. Vì ``quan hệ'' ở đây chỉ quan hệ lôgic giữa các
phần tử dữ liệu, cấu trúc dữ liệu còn được gọi là cấu trúc lôgic của dữ liệu.
Tương ứng với khái niệm khá trừu tượng ấy, cách biểu diễn cấu trúc dữ liệu
trong máy tính được gọi là cấu trúc lưu trữ của dữ liệu.

Lập mô hình toán học cho bài toán, rồi thiết kế thuật toán, viết chương trình
và gỡ lỗi cho đến khi chạy đúng, là các bước thông thường khi giải một bài
toán tin học. Muốn lập mô hình toán học, trước hết ta phải tìm ra quan hệ giữa
các đối tượng trong bài toán, tức xác định cấu trúc lôgic sẽ dùng. Đồng thời,
trong quá trình thiết kế thuật toán và hiện thực chương trình, ta phải xác
định cách thực hiện các thao tác trên từng đối tượng; cách thao tác lại phụ
thuộc vào cấu trúc lưu trữ được chọn. Do đó, cấu trúc lôgic và cấu trúc lưu
trữ tốt hay xấu sẽ trực tiếp ảnh hưởng tới hiệu suất chương trình.

\section{Chọn cấu trúc lôgic hợp lý}

Trong lập trình, chọn cấu trúc lôgic nghĩa là phân tích quan hệ giữa các phần
tử dữ liệu trong đề bài, rồi dựa trên những quan hệ đặc thù ấy để chọn cấu
trúc lôgic thích hợp nhằm mô tả bài toán bằng toán học và tiếp tục giải bài
toán. Thực chất, cấu trúc lôgic dùng phương pháp toán học để mô tả các đối
tượng thao tác trong bài toán và quan hệ giữa chúng: trừu tượng hóa đối tượng
thao tác thành phần tử toán học, và dùng ngôn ngữ toán học để mô tả các quan
hệ phức tạp giữa các đối tượng.

Theo đặc tính khác nhau của quan hệ giữa các phần tử dữ liệu, thường có bốn
cấu trúc lôgic cơ bản: tập hợp, cấu trúc tuyến tính, cấu trúc cây, và cấu trúc
đồ thị hay mạng. Trong bốn cấu trúc ấy, ngoại trừ tập hợp, nơi các phần tử chỉ
có quan hệ ``cùng thuộc một tập'', ba cấu trúc còn lại lần lượt biểu diễn các
quan hệ ``một-một'', ``một-nhiều'', và ``nhiều-nhiều''.

Vì vậy, trước khi chọn cấu trúc lôgic, ta phải phân tích rõ đối tượng thao tác
trong đề và quan hệ giữa các đối tượng ấy, rồi dựa vào đặc điểm của các quan
hệ để chọn cấu trúc lôgic hợp lý. Điều này đặc biệt quan trọng trong một số
bài toán phức tạp: quan hệ giữa dữ liệu rất rối, nhiều cấu trúc lôgic khác
nhau đều có thể giải bài toán, nhưng hiệu suất thuật toán thu được lại khác
nhau rất xa.

Với loại bài toán ấy, ta nên dùng tiêu chuẩn nào để chọn cấu trúc lôgic? Dưới
đây là hai yếu tố cần xét đầy đủ khi chọn cấu trúc lôgic hợp lý.

\subsection{Tận dụng thông tin có thể dùng trực tiếp}

Trước hết, ``thông tin'' ở đây là quan hệ giữa các phần tử.

Với thông tin cần xử lý, đại thể có thể chia thành hai loại: ``có thể dùng
trực tiếp'' và ``không thể dùng trực tiếp''. Với thông tin có thể dùng trực
tiếp, việc sử dụng rất thuận tiện: chỉ cần lấy ra là dùng được. Còn với loại
không thể dùng trực tiếp, ta cũng có thể thông qua một số cách gián tiếp để
biến nó thành thông tin dùng được, nhưng quá trình chuyển hóa ấy hiển nhiên
tốn thời gian.

Do đó, điều ta cần là càng nhiều thông tin có thể dùng trực tiếp càng tốt.
Loại thông tin này càng nhiều, hiệu suất thuật toán càng cao.

Các cấu trúc lôgic khác nhau chứa lượng thông tin khác nhau, và độ phức tạp
khi thuật toán sử dụng thông tin cũng khác nhau. Vì vậy, để thuật toán tận
dụng đầy đủ thông tin có thể dùng trực tiếp, đồng thời tránh lãng phí quá
nhiều thời gian trong quá trình chuyển thông tin không dùng trực tiếp thành
thông tin dùng được, ta cần dùng một cấu trúc lôgic hợp lý, sao cho nó chứa
nhiều thông tin có thể dùng trực tiếp hơn.

\paragraph{Bài toán 1: \textit{Hidden Code} của IOI 1999.}

\paragraph{Mô tả bài toán.}
Bài toán cho một số mã từ và một văn bản. Cần viết chương trình tìm mọi mục
mà văn bản chứa các mã từ ấy, rồi ghép các mục tìm được thành một ``đáp án''
tối ưu sao cho tổng độ dài các mã từ trong đáp án là lớn nhất. Mỗi mục gồm
một mã từ và một dãy phủ của mã từ ấy trong văn bản; chẳng hạn
\texttt{abcadc} là một dãy phủ của mã từ \texttt{abac}. Độ dài dãy phủ không
vượt quá \(1000\). Đồng thời, đáp án yêu cầu các dãy phủ của các mục không
chồng lấn nhau.

\paragraph{Phân tích.}
Với bài này, một thuật toán cơ bản khá dễ nghĩ ra là: lặp theo vị trí kết thúc
của dãy phủ trong văn bản, xét xem những mã từ nào được chứa, tìm tất cả các
mục, rồi cuối cùng dùng quy hoạch động để ghép các mục thành đáp án tối ưu.

Tạm thời ta chưa xét các mặt khác của thuật toán, mà trước hết chọn cấu trúc
lôgic cho bài toán.

Nếu dùng cấu trúc tuyến tính, chẳng hạn hàng đợi vòng, thì khi xét mã từ
\(t\), phương pháp là: ban đầu để con trỏ \(p\) trỏ tới vị trí kết thúc, sau
đó liên tục dịch \(p\) về trước để lần lượt tìm các chữ cái của mã từ \(t\)
từ cuối lên đầu. Ví dụ mã từ là \texttt{ABDCAB}; văn bản có thể hình dung như
sơ đồ sau, trong đó mũi tên ngoài cùng bên phải là vị trí kết thúc và mỗi mũi
tên ứng với một chữ cái được tìm thấy:

\begin{verbatim}
Huong di chuyen cua con tro p: tu phai sang trai

      A    B      D C       A        B

  C       D A C   B   D C   A D C        D B   A D C    C   B   A D

Hinh 1-1
\end{verbatim}

Vì đề quy định độ dài dãy phủ của mã từ không vượt quá \(1000\), mỗi lần xét
``có chứa hay không'' như vậy có độ phức tạp \(O(1000)\).

Trong văn bản, vị trí của hai chữ cái kề nhau trong mã từ \(t\) không nhất
thiết chỉ có quan hệ kề nhau; ví dụ \texttt{D} và \texttt{C} trong hình 1-1
có thể kề nhau, nhưng giữa \texttt{C} và \texttt{A} lại có thể cách hai chữ
cái. Thông tin trong cấu trúc tuyến tính chỉ tồn tại giữa hai phần tử kề
nhau. Dùng lượng thông tin đơn giản như vậy để tìm một chữ cái nào đó của mã
từ rõ ràng không hiệu quả.

Nếu ta dựng một đồ thị có hướng, trong đó đỉnh \(i\), tức vị trí thứ \(i\) của
văn bản, dùng \(52\) cung nối tới vị trí xuất hiện cuối cùng trước \(i\) của
từng chữ cái trong \texttt{a}..\texttt{z}, \texttt{A}..\texttt{Z}, thì khi cần
tìm chữ cái đứng trước một chữ cái nào đó trong mã từ, ta có thể dùng trực
tiếp cạnh đã nối, không cần liệt kê. Có thể xem bài toán là: từ một đỉnh của
đồ thị có hướng, tìm một đường đi độ dài \(\operatorname{length}(t)-1\), sao
cho các đỉnh đi qua có thứ tự theo các chữ cái trong mã từ \(t\).

\begin{verbatim}
  C       D A C   B   D C   A D C        D B   A D C    C   B   A D

Moi vi tri noi toi lan xuat hien gan nhat truoc do cua tung chu cai.

Hinh 1-2
\end{verbatim}

Tính toán cho thấy dùng đồ thị để ghi nhận là hoàn toàn chấp nhận được về
không gian: ghi \(1000\) điểm \(\times 52\) cung \(\times 4\) byte số nguyên
dài, tức khoảng \(200\) KB. Về thời gian, ta có thể tận dụng việc cách nối
cung của vị trí \(i\) và \(i+1\) thay đổi rất ít: ở hình 1-2, từ vị trí \(i\)
sang \(i+1\) chỉ có một cung đổi hướng, tức vị trí \(i+1\) cho một cung trỏ
tới vị trí \(i\). Vì vậy, muốn ghi các cung trong đồ thị, chỉ cần gán toàn bộ
hướng cung, rồi đổi một cung trong đó.

Nhờ dùng cấu trúc lôgic đồ thị, hiệu suất tìm chữ cái tăng mạnh. Độ phức tạp
của một lần xét là \(O(\operatorname{length}(t))\), tệ nhất là \(O(100)\),
tức nhanh hơn phương pháp ban đầu khoảng \(10\) lần. Chương trình kèm theo là
\texttt{codes.pas}.

Trong ví dụ này, tuy cấu trúc tuyến tính cũng có thể giải bài toán, nhưng do
tính đặc thù của phép kiểm tra, thông tin cần trong mỗi lần xét không thể tìm
từ các phần tử kề nhau. Việc cấu trúc tuyến tính chỉ có quan hệ giữa các phần
tử kề nhau trở thành một nhược điểm rõ rệt. Do đó, thông tin trong cấu trúc
tuyến tính của bài toán 1 thuộc loại không thể dùng trực tiếp. Ngược lại, cấu
trúc đồ thị vừa khớp với nhu cầu: nối bằng cung tất cả các điểm có thể phát
sinh quan hệ, để ta dùng quan hệ cung một cách hiệu quả trong quá trình xét
và tìm. Tuy cấu trúc đồ thị phức tạp hơn, nó đã chuyển thông tin không thể
dùng trực tiếp thành thông tin có thể dùng trực tiếp; hiệu suất thuật toán
tăng lên là điều tự nhiên.

\subsection{Không ghi nhận thông tin vô dụng}

Từ bài toán 1 ta thấy do cấu trúc đồ thị có lượng thông tin lớn, thông tin
trong đó về cơ bản đều ``có thể dùng''. Nhưng điều này không có nghĩa ta nhất
định phải dùng cấu trúc đồ thị. Trong một số tình huống, thông tin có thể dùng
trong cấu trúc đồ thị lại hơi thừa.

Thông tin đều dùng được đương nhiên là tốt, nhưng nếu trong đó có quá nhiều
thông tin vô dụng, tức không cần thiết, nó chỉ làm tăng độ phức tạp khi suy
nghĩ, phân tích và xử lý bài toán, rồi ngược lại gây bất lợi cho việc giải
bài.

\paragraph{Bài toán 2: Bài \textit{Đi thuyền} trong kỳ chọn đội Hồ Nam năm
1997.}

\paragraph{Mô tả bài toán.}
Có \(N\) người cần đi thuyền. Mỗi thuyền nhiều nhất chở hai người, và hai
người cùng thuyền phải cùng họ hoặc cùng tên. Hỏi cần ít nhất bao nhiêu
thuyền.

\paragraph{Phân tích.}
Nhìn bài này, nhiều người sẽ nghĩ tới cấu trúc dữ liệu đồ thị: xem \(N\)
người là \(N\) đỉnh của một đồ thị vô hướng, và nối cạnh giữa mọi cặp người
cùng họ hoặc cùng tên.

Điều kiện dùng ít thuyền nhất nghĩa là cần cho càng nhiều cặp người đi chung
một thuyền càng tốt; trong đồ thị, điều đó tương ứng với việc dùng ít cạnh
nhất để phủ tất cả các đỉnh. Đây vừa khớp với bài toán kinh điển trong lý
thuyết đồ thị: tìm phủ cạnh nhỏ nhất. Thuật toán tương ứng là thuật toán tìm
ghép cực đại trên đồ thị tổng quát.

Thuật toán tìm ghép cực đại trên đồ thị tổng quát khá phức tạp, cần dùng cây
luân phiên mở rộng, co hoa, v.v., và hiệu suất cũng không cao. Vì vậy, ta cần
tìm một phương pháp đơn giản và hiệu quả hơn.

Trước hết, do hai thành phần liên thông bất kỳ của đồ thị độc lập tương đối,
tức hai đỉnh của một cạnh ghép bất kỳ đều chỉ thuộc cùng một thành phần liên
thông, ta có thể xử lý riêng từng thành phần liên thông mà không ảnh hưởng
tới kết quả cuối cùng.

Đồng thời, ta có thể ước lượng cận dưới cho số thuyền \(s\). Với một thành
phần liên thông chứa \(P_i\) đỉnh, số thuyền tối thiểu hiển nhiên là
\(\lceil P_i/2\rceil\). Nếu đồ thị có tổng cộng \(L\) thành phần liên thông,
thì
\[
  s=\sum_{i=1}^{L}\left\lceil\frac{P_i}{2}\right\rceil .
\]
Sau nhiều lần thử, ta có thể đưa ra phỏng đoán: số cạnh phủ thực tế cần dùng
hoàn toàn bằng cận dưới vừa tính.

Nếu dùng cấu trúc đồ thị để chứng minh phỏng đoán trên, có thể dựa vào hai
bước sau:
\begin{enumerate}
  \item Nếu trong một thành phần liên thông không tồn tại đỉnh bậc \(1\), thì
  chắc chắn tồn tại chu trình.
  \item Xóa một đỉnh bậc \(1\) và đỉnh kề của nó, hoặc xóa bất kỳ cạnh nào
  trong một chu trình, thành phần liên thông vẫn liên thông; tức thành phần
  liên thông chắc chắn tồn tại cạnh không phải cầu.
\end{enumerate}
Vì phương pháp đồ thị không phải trọng tâm ở đây, chi tiết chứng minh không
trình bày thêm. Từ cấu trúc dữ liệu đồ thị, thuật toán thu được là: mỗi lần
in ra một cạnh không phải cầu, rồi xóa hai đỉnh của cạnh ấy khỏi đồ thị. Độ
phức tạp thời gian là \(O(n^3)\): tìm một cạnh không phải cầu mất \(O(n^2)\),
và thao tác tìm cạnh phủ mất \(O(n)\).

Do bị giới hạn bởi cấu trúc đồ thị, độ phức tạp thời gian đã khó giảm tiếp.
Vì vậy, nếu muốn tiếp tục tối ưu thuật toán, ta chỉ có thể cân nhắc dùng một
cấu trúc lôgic khác. Ở đây ta nghĩ tới cấu trúc cây nhị phân: cụ thể là
chuyển từng thành phần liên thông trong đồ thị thành cây nhị phân, rồi dùng
cây nhị phân để giải bài toán.

Trước hết chọn một đỉnh bất kỳ trong thành phần liên thông làm gốc, rồi xác
định cách dựng cây:
\begin{enumerate}
  \item Tìm một điểm \(j\) cùng họ với nút gốc \(i\), với \(j\) chưa thuộc
  cây nhị phân, làm con trái của \(i\), rồi lấy \(j\) làm gốc để dựng cây con.
  \item Tìm một điểm \(k\) cùng tên với nút gốc \(i\), với \(k\) chưa thuộc
  cây nhị phân, làm con phải của \(i\), rồi lấy \(k\) làm gốc để dựng cây con.
\end{enumerate}

Với một thành phần liên thông như hình 2-1-1, dùng cách dựng cây trên có thể
nhận được cây nhị phân như hình 2-1-2, lấy nút \(1\) làm gốc. Trong hình gốc,
đường liền biểu diễn cùng họ, đường đứt biểu diễn cùng tên. Vì hình chỉ là
sơ đồ nhỏ, có thể mô tả cây thu được như sau:
\[
  1\to(2,5),\qquad 2\to(3,\varnothing),\qquad
  3\to(4,\varnothing),\qquad 5\to(6,\varnothing),\qquad
  6\to(7,8).
\]

Tiếp theo ta chứng minh cây này chắc chắn chứa mọi đỉnh của thành phần liên
thông.

\paragraph{Bổ đề 2.1.}
Nếu cây nhị phân \(T\) chứa một nút \(p\), thì tất cả các điểm cùng họ với
\(p\) trong thành phần liên thông đều thuộc \(T\).

\paragraph{Chứng minh.}
Để tiện trình bày, quy ước \(S\) là tập đỉnh cùng họ với \(p\);
\(\operatorname{lc}[p,0]\) là nút \(p\);
\(\operatorname{lc}[p,i]\), với \(i>0\), là con trái của
\(\operatorname{lc}[p,i-1]\). Hiển nhiên \(\operatorname{lc}[p,i]\) cùng họ
với \(p\).

Giả sử tồn tại một điểm \(q\) thỏa \(q\in S\) nhưng \(q\notin T\). Vì \(S\) là
tập hữu hạn, chắc chắn tồn tại một \(\operatorname{lc}[p,k]\) không có con
trái. Khi đó ta có thể đặt \(\operatorname{lc}[p,k+1]=q\), suy ra \(q\in T\),
mâu thuẫn với giả thiết \(q\notin T\). Vậy giả thiết sai và bổ đề được chứng
minh.

Từ cách chứng minh bổ đề 2.1, tương tự ta có bổ đề 2.2.

\paragraph{Bổ đề 2.2.}
Nếu cây nhị phân \(T\) chứa một nút \(p\), thì tất cả các điểm cùng tên với
\(p\) trong thành phần liên thông đều thuộc \(T\).

\paragraph{Định lý 1.}
Cây nhị phân lấy bất kỳ điểm \(p\) nào trong một thành phần liên thông làm
gốc chắc chắn chứa mọi đỉnh của thành phần liên thông ấy.

\paragraph{Chứng minh.}
Theo bổ đề 2.1 và bổ đề 2.2, tất cả các điểm cùng họ hoặc cùng tên với \(p\)
đều chắc chắn ở trong cây nhị phân, tức mọi điểm có cạnh nối với \(p\) trong
thành phần liên thông đều ở trong cây. Do giữa hai điểm bất kỳ trong một
thành phần liên thông đều tồn tại đường đi, mọi điểm của thành phần liên
thông ấy đều thuộc cây nhị phân.

Sau khi chứng minh cây nhị phân chứa mọi đỉnh của thành phần liên thông, ta
cần chứng minh phỏng đoán ban đầu, tức định lý sau.

\paragraph{Định lý 2.}
Với cây nhị phân \(T_m\) chứa \(m\) nút, số thuyền cần dùng là
\[
  \operatorname{boat}[m]=\left\lceil\frac{m}{2}\right\rceil,\qquad m\in N.
\]

\paragraph{Chứng minh.}
\begin{enumerate}[label=(\roman*)]
  \item Khi \(m=1,2,3\), mệnh đề hiển nhiên đúng.
  \item Giả sử mệnh đề đúng với \(m<k\), \(k>3\). Khi \(m=k\), trước hết tìm
  một nút có tầng sâu nhất trong cây, và gọi cha của nút ấy là \(p\). Khi đó
  chỉ có ba tình huống sau.
\end{enumerate}

\begin{enumerate}
  \item \(p\) chỉ có một con. Khi đó xóa \(p\) và người con duy nhất của
  \(p\), \(T_k\) trở thành \(T_{k-2}\), nên
  \[
    \operatorname{boat}[k]=\operatorname{boat}[k-2]+1
    =\left\lceil\frac{k-2}{2}\right\rceil+1
    =\left\lceil\frac{k}{2}\right\rceil .
  \]
  \item \(p\) có hai con, và \(p\) là con trái của cha nó. Khi đó có thể xóa
  \(p\) và con phải của \(p\), rồi đặt con trái của \(p\) vào vị trí của
  \(p\). Tương tự, \(T_k\) trở thành \(T_{k-2}\), nên
  \(\operatorname{boat}[k]=\operatorname{boat}[k-2]+1
  =\lceil k/2\rceil\).
  \item \(p\) có hai con, và \(p\) là con phải của cha nó. Khi đó có thể xóa
  \(p\) và con trái của \(p\), rồi đặt con phải của \(p\) vào vị trí của
  \(p\). Tình huống này rất giống trường hợp 2, nên cũng có
  \(\operatorname{boat}[k]=\operatorname{boat}[k-2]+1=\lceil k/2\rceil\).
\end{enumerate}
Kết hợp ba trường hợp, khi \(m=k\) ta có
\(\operatorname{boat}[k]=\lceil k/2\rceil\). Cuối cùng, từ (i) và (ii), với
mọi \(m\in N\), \(\operatorname{boat}[m]=\lceil m/2\rceil\).

Từ chứng minh trên, sau khi chuyển cấu trúc đồ thị của dữ liệu thành cấu trúc
cây, ta có thuật toán tìm phương án đi thuyền cho một cây nhị phân:

\begin{verbatim}
proc try(father: integer; var root: integer; var rest: byte);
{In phuong an di thuyen trong cay con goc root.
 father=0 nghia la root la con trai cua cha no,
 father=1 nghia la root la con phai cua cha no,
 rest cho biet sau khi in phuong an cho cay con co con lai
 mot nut goc chua di thuyen hay khong.}
begin
  visit[root] := true; {danh dau root da tham}
  tim mot nut j cung ho voi root va chua tham;
  if j <> n+1 then try(0, j, lrest);
  tim mot nut k cung ten voi root va chua tham;
  if k <> n+1 then try(1, k, rrest);
  if (lrest=1) xor (rrest=1) then begin
    {root chi co mot con, truong hop 1}
    if lrest=1 then print(lrest, root) else print(rrest, root);
    rest := 0;
  end
  else if (lrest=1) and (rrest=1) then begin
    {root co hai con}
    if father=0 then begin
      print(rrest, root); root := j; {truong hop 2}
    end
    else begin
      print(lrest, root); root := k; {truong hop 3}
    end;
    rest := 1;
  end
  else rest := 1;
end;
\end{verbatim}

Đây chỉ là thuật toán in phương án đi thuyền cho một cây nhị phân. Muốn in
phương án cho tất cả mọi người, ta cần thêm một vòng lặp để tìm gốc của từng
cây nhị phân. Vì mỗi điểm chỉ được thăm một lần, và mỗi lần tìm con trái,
con phải đều cần một vòng lặp, độ phức tạp thời gian của thuật toán là
\(O(n^2)\). Chương trình kèm theo là \texttt{boat.pas}.

Cuối cùng, ta phân tích nguyên nhân hai cấu trúc dẫn tới hai thuật toán có độ
phức tạp khác nhau. Điểm mấu chốt là: tuy cây nhị phân có cấu trúc tương đối
đơn giản, nó đã chứa gần như toàn bộ thông tin hữu dụng. Từ thuật toán tìm
phương án đi thuyền, ta thấy mọi cạnh trong cây nhị phân không chỉ đều phát
huy tác dụng, mà còn không bị dùng lặp lại; tỷ lệ sử dụng thông tin khá cao.

Khi cấu trúc cây đã đủ, một số thông tin trong cấu trúc đồ thị hiển nhiên trở
thành thông tin vô dụng. Những thông tin vô dụng dư thừa ấy làm ta khó phát
hiện quy luật khi phân tích bài toán, và cũng khó tìm thuật toán hiệu quả.
Cũng như tường trong mê cung, càng nhiều càng khó đi. Thông tin vô dụng chỉ
làm nhiễu tính quy luật của bài toán, khiến ta khó tìm phương pháp giải hơn.

\paragraph{Tiểu kết.}
Chọn cấu trúc lôgic của dữ liệu là một khâu then chốt trong việc xây dựng mô
hình toán học, còn thuật toán dùng để giải mô hình toán học ấy. Muốn thuật
toán hiệu quả, trước hết phải chọn tốt cấu trúc lôgic của dữ liệu. Trên đây
đã nêu hai điều kiện, hay hai hướng suy nghĩ, khi chọn cấu trúc lôgic; mục
đích chung là nâng cao hiệu quả sử dụng thông tin. Thông tin có thể dùng trực
tiếp không cần thao tác trung gian nên hiệu quả sử dụng tự nhiên cao. Không
ghi nhận thông tin vô dụng giúp ta tập trung hơn vào việc nghiên cứu và phân
tích thông tin hữu dụng, và việc sử dụng thông tin cũng nhờ đó được tối ưu
hơn.

Tóm lại, trong quá trình giải bài toán, chọn cấu trúc lôgic hợp lý là một
khâu rất quan trọng.

\section{Chọn cấu trúc lưu trữ hợp lý}

Cấu trúc lưu trữ của dữ liệu được chia thành cấu trúc lưu trữ tuần tự và cấu
trúc lưu trữ liên kết. Đặc điểm của cấu trúc lưu trữ tuần tự là dùng vị trí
tương đối của phần tử trong bộ nhớ để biểu diễn quan hệ lôgic giữa các phần
tử dữ liệu. Cấu trúc lưu trữ liên kết thì dùng con trỏ chỉ tới địa chỉ lưu
trữ của phần tử để biểu diễn quan hệ lôgic giữa các phần tử dữ liệu.

Do hai cấu trúc lưu trữ khác nhau, khi sử dụng cụ thể chúng cũng có ưu điểm
và nhược điểm riêng.

Xét một ví dụ đơn giản: cần ghi một ma trận \(n\times n\), trong đó có \(m\)
phần tử khác \(0\). Nếu dùng cấu trúc lưu trữ tuần tự, ta sẽ dùng một mảng
hai chiều \(n\times n\), ghi lại toàn bộ phần tử dữ liệu. Nếu dùng cấu trúc
lưu trữ liên kết, ta cần một danh sách liên kết gồm \(m\) nút, ghi lại \(m\)
phần tử dữ liệu khác \(0\). Từ hai cách ghi khác nhau ấy, ta có thể phân tích
ưu nhược điểm của chúng qua các thao tác khác nhau trên dữ liệu.

\begin{enumerate}
  \item Truy cập ngẫu nhiên một phần tử bất kỳ trong ma trận. Do cấu trúc tuần
  tự liền nhau về vị trí vật lý, rất dễ lấy địa chỉ lưu trữ của bất kỳ phần
  tử nào, độ phức tạp là \(O(1)\). Với cấu trúc liên kết, vì không có tính
  liền nhau về vị trí vật lý, trước hết phải duyệt toàn bộ danh sách để tìm
  địa chỉ lưu trữ của phần tử cần truy cập, độ phức tạp là \(O(m)\). Lúc này
  dùng cấu trúc tuần tự rõ ràng hiệu quả hơn.
  \item Duyệt toàn bộ dữ liệu. Độ phức tạp của hai cấu trúc lưu trữ đối với
  thao tác này rất rõ: cấu trúc tuần tự là \(O(n^2)\), cấu trúc liên kết là
  \(O(m)\). Vì thông thường \(m\) nhỏ hơn \(n^2\) rất nhiều, lúc này cấu trúc
  liên kết hiệu quả hơn nhiều.
\end{enumerate}

Ngoài hai thao tác trên, với các thao tác khác, hai cấu trúc này không có ưu
nhược điểm quá rõ rệt. Chẳng hạn, xóa hoặc chèn trên danh sách liên kết có
thể được biểu diễn trong cấu trúc tuần tự bằng cách đổi phần tử dữ liệu ở vị
trí tương ứng.

Vì hiệu suất của hai cấu trúc lưu trữ đối với các thao tác khác nhau chênh
lệch khá lớn, khi xác định cấu trúc lưu trữ, ta phải phân tích cẩn thận nhu
cầu thao tác trong thuật toán, rồi chọn hợp lý một cấu trúc có thể phát huy
điểm mạnh và tránh điểm yếu.

\subsection{Dùng cấu trúc lưu trữ tuần tự một cách hợp lý}

Khi làm bài hằng ngày, phần lớn chúng ta dùng cấu trúc lưu trữ tuần tự để lưu
dữ liệu. Nguyên nhân một mặt là thao tác trên cấu trúc tuần tự thuận tiện,
mặt khác trong quá trình hiện thực chương trình, dùng cấu trúc tuần tự tiện
gỡ lỗi và tìm lỗi hơn so với cấu trúc liên kết. Vì vậy, theo thói quen, đa số
người cho rằng bài toán nào có thể dùng cấu trúc tuần tự để lưu trữ thì tốt
nhất nên dùng cấu trúc lưu trữ tuần tự.

Thực ra, cái ``tốt'' ấy chỉ là một tiêu chuẩn tương đối, được xây dựng trên
hai điều kiện tiên quyết sau:
\begin{enumerate}
  \item Số nút được lưu bằng cấu trúc liên kết không khác nhiều so với số nút
  được lưu bằng cấu trúc tuần tự. Trong trường hợp này, vì số nút cần lưu khá
  gần nhau, dùng cấu trúc liên kết hoàn toàn không thể thể hiện ưu điểm ghi ít
  nút hơn, và thậm chí còn có thể làm giảm hiệu suất thuật toán do thao tác
  con trỏ chậm hơn. Nghiêm trọng hơn, vì bản thân con trỏ chiếm nhiều không
  gian, lại có nhiều nút, yêu cầu không gian của thuật toán có thể không được
  đáp ứng.
  \item Đây không phải nút thắt hiệu suất của thuật toán. Vì không phải phần
  tốn thời gian nhất của thuật toán, việc cải tiến ở đây hiển nhiên không ảnh
  hưởng lớn tới toàn bộ thuật toán; nếu dùng cấu trúc liên kết, thao tác còn
  trở nên rườm rà quá mức.
\end{enumerate}

\subsection{Dùng cấu trúc lưu trữ liên kết khi cần thiết}

Trên đây đã phân tích điều kiện dùng cấu trúc lưu trữ tuần tự. Vấn đề còn lại
là khi nào nên dùng cấu trúc lưu trữ liên kết.

Do thao tác con trỏ trong cấu trúc liên kết quả thật khá rườm rà, tốc độ cũng
chậm hơn và gỡ lỗi không thuận tiện, mọi người thường không muốn dùng cấu
trúc lưu trữ liên kết. Nhưng đó chỉ là quan điểm thông thường; khi cấu trúc
liên kết thật sự cải thiện thuật toán rất lớn, ta vẫn buộc phải cân nhắc.

\paragraph{Bài toán 3: \textit{Underground City} của IOI 1999.}

\paragraph{Mô tả bài toán.}
Đã biết bản đồ của một thành phố, nhưng không cho biết vị trí ban đầu của
bạn. Bạn cần thông qua một chuỗi bước di chuyển và thăm dò để xác định vị trí
ban đầu. Các ràng buộc của đề là:
\begin{enumerate}
  \item Không được di chuyển tới ô có tường.
  \item Chỉ được thăm dò ô kề với vị trí hiện tại theo bốn hướng.
\end{enumerate}
Dưới hai ràng buộc này, số lần thăm dò, không tính di chuyển, cần càng ít
càng tốt.

\paragraph{Phân tích.}
Vì cấu trúc lưu trữ phải do nhu cầu của thuật toán quyết định, trước hết ta
xác định thuật toán cho bài toán.

Sau khi phân tích bài, ta có tư tưởng cơ bản: ban đầu giả sử mọi ô không có
tường đều có thể là vị trí xuất phát, rồi thông qua thăm dò để từng bước thu
hẹp phạm vi vị trí xuất phát, cuối cùng nhận được vị trí xuất phát thật. Đồng
thời, để nâng hiệu suất thuật toán, ta còn dùng tư tưởng chia để trị, sao cho
mỗi lần thăm dò đều cố gắng thu hẹp phạm vi vị trí xuất phát nhiều nhất có
thể, tức để chương trình phụ thuộc vào may mắn ít nhất.

Tiếp đó, ta xác định cấu trúc lưu trữ cho bài này.

Vì bản đồ trong bài là một ma trận hai chiều, thông thường dùng cấu trúc lưu
trữ tuần tự là điều tự nhiên. Nhưng trong bài này, cấu trúc lưu trữ tuần tự
bộc lộ nhược điểm lớn. Các thao tác ta thực hiện nhiều nhất là lọc phạm vi vị
trí xuất phát và quyết định chọn vị trí nào để thăm dò. Dữ liệu cần dùng cho
hai thao tác ấy chỉ là một phần rất nhỏ của bản đồ lớn. Nếu dùng cấu trúc lưu
trữ tuần tự, dù cần dùng bao nhiêu dữ liệu, ta vẫn luôn phải duyệt toàn bộ
bản đồ.

\begin{verbatim}
Hinh 3-1: dung mang hai chieu; phan da danh dau chiem it o,
nhung van phai quet ca ban do.

Hinh 3-2: dung danh sach lien ket:
head -> (2,2) -> (2,3) -> (3,4) -> (4,1) -> nil
\end{verbatim}

Nếu dùng cấu trúc lưu trữ liên kết như danh sách ở hình 3-2, cần bao nhiêu dữ
liệu thì chỉ duyệt bấy nhiêu dữ liệu. Cách này không chỉ phát huy đầy đủ ưu
điểm của cấu trúc liên kết, mà vì không cần lấy riêng một dữ liệu nào, mỗi
lần đều xét toàn bộ dữ liệu trong danh sách, nó cũng tránh được nhược điểm
lớn nhất của cấu trúc liên kết.

Dùng cấu trúc lưu trữ liên kết tuy không làm giảm độ phức tạp thời gian của
bài toán, vì trong trường hợp xấu nhất lượng lưu trữ của cấu trúc liên kết
gần như bằng cấu trúc tuần tự, nhưng do thể hiện nguyên tắc phát huy điểm
mạnh và tránh điểm yếu khi chọn cấu trúc lưu trữ, hiệu suất thuật toán vẫn
tăng rất lớn. Thời gian chạy trên các bộ dữ liệu khác nhau như bảng 3-3.

\begin{center}
\begin{tabular}{c r r}
\hline
Bộ dữ liệu & Cấu trúc tuần tự & Cấu trúc liên kết\\
\hline
1 & 0.06s & 0.02s\\
2 & 1.73s & 0.07s\\
3 & 1.14s & 0.06s\\
4 & 3.86s & 0.14s\\
5 & 32.84s & 0.21s\\
6 & 141.16s & 0.23s\\
7 & 0.91s & 0.12s\\
8 & 6.92s & 0.29s\\
9 & 6.10s & 0.23s\\
10 & 17.41s & 0.20s\\
\hline
\end{tabular}
\end{center}

Chương trình dùng cấu trúc lưu trữ liên kết kèm theo là \texttt{under.pas}.

Ta chọn cấu trúc lưu trữ liên kết tuy thao tác có thể phức tạp hơn đôi chút,
nhưng vì cải tiến đúng nút thắt của thuật toán, hiệu suất thuật toán đã khác
hẳn. Qua đó có thể thấy, hiệu quả của việc chọn cấu trúc liên kết khi cần
thiết là không thể xem nhẹ.

\paragraph{Tiểu kết.}
Chọn cấu trúc lôgic hợp lý liên quan tới việc xây dựng mô hình toán học, nên
có lẽ mọi người đều chú ý hơn. Nhưng việc chọn cấu trúc lưu trữ lại dễ bị bỏ
qua vì không làm thay đổi độ phức tạp của thuật toán. Vậy một phương pháp
không làm giảm độ phức tạp thuật toán có đáng coi trọng không?

Ai cũng biết, cắt tỉa là một phương pháp tối ưu hóa thuật toán thường dùng và
được sử dụng rộng rãi, nhưng cắt tỉa cũng không thể thay đổi độ phức tạp của
thuật toán. Vì vậy, việc chọn hợp lý cấu trúc lưu trữ, có tác dụng tương tự
cắt tỉa, cũng rất đáng coi trọng.

Tóm lại, trong quá trình thiết kế thuật toán, ta phải xét đầy đủ các ảnh
hưởng khác nhau do cấu trúc lưu trữ mang lại, rồi chọn cấu trúc lưu trữ hợp
lý nhất.

\section{Kết hợp nhiều cấu trúc dữ liệu}

Những phần trên đều thảo luận cách chọn cấu trúc dữ liệu, gồm lựa chọn cấu
trúc lôgic và lựa chọn cấu trúc lưu trữ; đây là một phương pháp tối ưu hóa
thuật toán có tính phổ biến khá lớn. Với phần lớn bài toán, ta đều có thể
thông qua việc chọn một cấu trúc lôgic và một cấu trúc lưu trữ hợp lý để đạt
mục đích tối ưu thuật toán.

Nhưng một số bài toán lại thường không thuận như vậy. Khi chọn cấu trúc dữ
liệu cho loại bài toán này, ta thường được cái này mất cái kia; thậm chí có
khi hoàn toàn không tồn tại một cấu trúc dữ liệu thích hợp. Lúc ấy ta không
thể chọn ra một cấu trúc dữ liệu đơn lẻ phù hợp, và các phương pháp trên
không còn thích hợp lắm.

Để giải quyết khó khăn trong việc chọn cấu trúc dữ liệu, ta có thể dùng
phương pháp kết hợp nhiều cấu trúc dữ liệu. Thông qua kết hợp nhiều cấu trúc
dữ liệu, ta đạt được tác dụng bù trừ ưu nhược điểm, để các cấu trúc dữ liệu
khác nhau phát huy ưu thế riêng trong thuật toán.

Đó chỉ là tư tưởng chung khi kết hợp nhiều cấu trúc dữ liệu. Cụ thể kết hợp
ra sao, ta xem ví dụ sau.

\paragraph{Bài toán 4: Bài \textit{Dãy nhỏ nhất} trong kỳ thi tỉnh Quảng Đông
năm 1997.}

\paragraph{Mô tả bài toán.}
Cho một ma trận số nguyên dương \(N\times N\), với \(N\le 100\). Cần tìm một
đường đi từ vị trí đầu đến vị trí cuối trong ma trận, có thể đi theo bốn
hướng lên, xuống, trái, phải, sao cho tổng trị tuyệt đối của hiệu giữa hai số
kề nhau trên đường đi là nhỏ nhất.

\paragraph{Phân tích.}
Thuật toán cơ bản của bài này rất đơn giản: chỉ cần dùng thuật toán Dijkstra
để tìm đường đi ngắn nhất từ vị trí đầu đến vị trí cuối. Nhưng ở đây có một
vấn đề lớn: \(N\le 100\), tức số điểm trong đồ thị có thể lên tới \(10000\).
Khi đó thuật toán Dijkstra độ phức tạp \(O(n^2)\) tỏ ra hơi quá sức.

Ta tiếp tục phân tích thuật toán. Do Dijkstra thường dùng cấu trúc mảng tuyến
tính, mỗi lần tìm đường đi ngắn tiếp theo có hai bước:
\begin{enumerate}
  \item Tìm một đỉnh \(i\) không thuộc tập đỉnh xuất phát của các đường đi
  ngắn nhất, và có khoảng cách tới đích nhỏ nhất. Độ phức tạp của bước này
  hiển nhiên là \(O(n)\).
  \item Sửa độ dài đường đi từ các đỉnh kề với \(i\) tới đích. Vì nhiều nhất
  chỉ có \(4\) điểm, tương ứng bốn hướng, kề với \(i\), độ phức tạp của bước
  này là \(O(1)\), tức hằng số.
\end{enumerate}

Trong hai bước này, tuy bước thứ hai nhiều nhất chỉ đổi độ dài đường đi tới
\(4\) điểm, bước thứ nhất vẫn phải liệt kê tất cả phần tử của mảng, rõ ràng
rất tốn thời gian. Để cải tiến bước thứ nhất, ta có thể nghĩ tới cấu trúc
đống nhị phân trong cấu trúc cây. Ưu điểm của đống là: nút gốc chính là nút
tối ưu, và sau khi đổi giá trị của một nút trong đống, chỉ cần độ phức tạp
\(O(\log_2 n)\) để hoàn tất quá trình vun đống. Quá trình này có thể chia
thành vun từ dưới lên hoặc từ trên xuống trong cây nhị phân, và độ phức tạp
đều như nhau. Nhờ đó, sau khi cấu trúc đống thay đổi, vun đống vẫn giữ được
tính chất của đống.

Nếu dùng cấu trúc đống nhị phân để lưu khoảng cách, bước thứ nhất chỉ cần lấy
nút gốc ra, dùng một nút khác trong đống thay thế rồi vun từ trên xuống một
lần. Vì thế độ phức tạp của bước thứ nhất có thể giảm xuống \(O(\log_2 n)\).
Nhưng lúc này cấu trúc đống bộc lộ nhược điểm lớn ở bước thứ hai: ta không
biết trước vị trí của các điểm kề với \(i\) trong đống. Nếu duyệt đống để tìm
các điểm kề với \(i\), độ phức tạp của bước thứ hai lại trở thành \(O(n)\).

Từ hai cấu trúc dữ liệu này, không khó thấy quy luật: dùng mảng tuyến tính
thì dễ thực hiện bước thứ hai; dùng đống nhị phân của cấu trúc cây thì bước
thứ nhất có hiệu quả lý tưởng hơn.

Vậy ta có thể kết hợp hai cấu trúc dữ liệu ấy, lấy dài bù ngắn không? Câu trả
lời là có.

Ta có thể dùng phương pháp ánh xạ để thiết lập tương ứng một-một giữa phần tử
trong cấu trúc tuyến tính và nút trong đống. Nếu phần tử trong mảng tuyến
tính thay đổi, nút tương ứng trong đống cũng thay đổi; nếu nút trong đống
thay đổi, phần tử tương ứng trong mảng cũng đổi theo.

Sau khi kết hợp hai cấu trúc, dù là bước thứ nhất hay thứ hai, ta đều không
cần duyệt tất cả phần tử, mà chỉ cần thực hiện một số hằng định lần thao tác
vun đống độ phức tạp \(O(\log_2 n)\). Như vậy, độ phức tạp thời gian toàn bộ
trở thành \(O(n\log_2 n)\), hiệu suất thuật toán rõ ràng được nâng cao rất
nhiều.

Ví dụ này chỉ phân tích sự kết hợp của cấu trúc lôgic. Với cấu trúc lưu trữ,
thực ra cũng cùng một đạo lý. Ta cũng có thể thiết lập quan hệ tương ứng giữa
cấu trúc lưu trữ liên kết và cấu trúc lưu trữ tuần tự, đồng thời thực hiện
song song thao tác chèn, xóa trong cấu trúc liên kết với việc thay đổi dấu
hiệu phần tử dữ liệu trong cấu trúc tuần tự. Như vậy, khi truy cập phần tử dữ
liệu, nếu cần duyệt thì dùng cấu trúc liên kết; nếu chỉ cần xét một phần tử
nào đó, thì truy cập trực tiếp dấu hiệu của nó trong cấu trúc tuần tự. Thông
qua kết hợp cấu trúc lưu trữ, ta phát huy được đầy đủ ưu điểm của cả hai cấu
trúc.

Từ đó, ta lại có thêm một phương pháp chọn và sử dụng cấu trúc dữ liệu khi ưu
khuyết điểm của các cấu trúc khó phân biệt và khó lựa chọn: dùng ánh xạ để
kết hợp các cấu trúc dữ liệu. Đây rõ ràng là một phương pháp và kỹ xảo nữa
trong việc sử dụng cấu trúc dữ liệu.

\section{Kết luận}

Khi dùng cấu trúc dữ liệu hằng ngày, ta thường chỉ xem nó là công cụ để xây
dựng mô hình và hiện thực thuật toán, mà chưa xét ảnh hưởng của công cụ ấy
đối với hiệu suất chương trình. Khi độ khó của các bài toán tin học ngày càng
tăng, yêu cầu về hiệu suất thời gian và không gian của thuật toán cũng ngày
càng cao; mà hiệu suất thời gian và không gian của thuật toán, ở mức rất lớn,
đều bị cấu trúc dữ liệu chi phối.

Vì vậy, cấu trúc dữ liệu là nền tảng của thi tin học và cũng là nội dung mọi
người quen thuộc nhất, nhưng nó giữ vai trò then chốt đối với cả thời gian
lẫn không gian của thuật toán. Chỉ khi phân tích ở mức nguyên tắc các phương
pháp dùng cấu trúc dữ liệu để tối ưu thuật toán, và khi làm bài xét đầy đủ
nguyên tắc lựa chọn phát huy điểm mạnh, tránh điểm yếu cùng phương pháp kết
hợp lấy dài bù ngắn, ta mới có thể cải tiến thuật toán hiệu quả hơn.

Cùng với sự phát triển không ngừng của cấu trúc dữ liệu, nhận thức của con
người về cấu trúc dữ liệu cũng ngày càng sâu sắc. Bài viết này chỉ thảo luận
sơ bộ về cấu trúc dữ liệu; muốn hiểu và nắm vững cấu trúc dữ liệu một cách
toàn diện, ta còn cần phân tích và tổng kết sâu hơn. Như câu nói ``nuôi quân
nghìn ngày, dùng quân một buổi'', chỉ khi ngày thường xây nền tảng tốt và
tổng kết kinh nghiệm, ta mới có thể vận dụng thành thạo trong các kỳ thi tin
học.

\appendix

\section{Phụ chú}

\begin{enumerate}
  \item Định nghĩa cấu trúc dữ liệu trong bài viết này dùng định nghĩa trong
  sách \textit{Cấu trúc dữ liệu, bản thứ hai}; trong các sách khác, định
  nghĩa cấu trúc dữ liệu có thể khác.
  \item Vì trọng tâm bài viết là cấu trúc dữ liệu chứ không phải thuật toán,
  khi phân tích bài toán, bài viết có thể chưa trình bày thật chi tiết quá
  trình suy ra thuật toán của từng bài.
  \item Khi so sánh chương trình dùng cấu trúc lưu trữ tuần tự và cấu trúc
  lưu trữ liên kết trong bài \textit{Underground City}, dữ liệu được chọn là
  dữ liệu chuẩn của IOI 1999. Máy thử nghiệm là 6x86 200MHz.
\end{enumerate}

\section{Danh sách chương trình}

Vì bài viết thảo luận phương pháp chọn cấu trúc dữ liệu, khi phân tích ví dụ
phần lớn dùng cách so sánh các cấu trúc dữ liệu khác nhau. Dưới đây chỉ đưa
ra các chương trình dùng cấu trúc dữ liệu tốt hơn.

\subsection{\texttt{codes.pas}: \textit{Hidden Code}, dùng cấu trúc đồ thị có hướng}

Để nâng hiệu suất chương trình, trong chương trình còn có một số biện pháp
tối ưu khác, như sắp xếp mã từ và lưu chọn lọc các mục tìm được. Định dạng
vào ra theo chuẩn của đề thi.

\begin{verbatim}
{$A+,B-,D+,E+,F-,G-,I-,L+,N-,O-,P-,Q-,R-,S-,T-,V+,X+,Y+}
{$M 2048,0,655360}
const st11='words.inp';   {ten tep ghi cac ma tu}
      st12='text.inp';    {ten tep ghi van ban}
      st2='codes.out';    {ten tep xuat}
      max=1023;           {so chu cai toi da ghi bang hang doi vong}
type arr1=array[0..10000] of longint;
     arr2=array[0..10000] of shortint;
     arr3=array['A'..'z'] of longint;
var i,j,now:longint; {now danh dau da doc den vi tri thu now cua van ban}
    m,n:integer;     {n la so ma tu, m la so muc tim duoc}
    buf:array[1..4095] of char; {bo dem tep vao}
    t:array[1..100] of string[100]; {t[i] ghi tung ma tu}
    last:array[0..max] of ^arr3;
    {vi tri xuat hien cuoi cung cua chu c truoc vi tri i duoc ghi
     la last[i and max,c]}
    nearest:arr3; {nearest[c] la vi tri xuat hien cuoi cung cua c
                   truoc vi tri now hien tai}
    bb,ee:array['A'..'z'] of byte;
    {sau khi sap xep, cac ma tu co chu cuoi c nam o bb[c]..ee[c];
     bb[c]>ee[c] nghia la khong co ma tu co chu cuoi c}
    no:array[1..100] of byte;
    {no[i] la vi tri trong danh sach goc cua ma tu thu i sau sap xep}
    ss,tt:^arr1;
    nn:^arr2;
    {ss[i],tt[i],nn[i] lan luot ghi vi tri dau, vi tri cuoi cua day phu
     trong muc thu i va so hieu ma tu tuong ung trong danh sach da sap xep}
    f:text;
    b:char;

procedure readcod; {doc danh sach ma tu va sap xep}
var f:text;
    i,j,k,l:integer;
    st:string;
begin
   assign(f,st11); reset(f);
   readln(f,n);
   for i:=1 to n do begin
       readln(f,t[i]); no[i]:=i;
   end;
   fillchar(bb,sizeof(bb),1);
   fillchar(ee,sizeof(ee),0);
   {ban dau de moi bb[c]>ee[c], sau khi tim chu c moi sua bb[c], ee[c]}
   for i:=1 to n do begin
       k:=i;
       for j:=i+1 to n do
           {sap theo chu cuoi; neu chu cuoi bang nhau thi ma tu dai hon truoc}
           if (t[j,length(t[j])]<t[k,length(t[k])]) or
              (t[j,length(t[j])]=t[k,length(t[k])]) and
              (length(t[j])>length(t[k])) then k:=j;
       if i<>k then begin
          st:=t[i]; t[i]:=t[k]; t[k]:=st;
          l:=no[i]; no[i]:=no[k]; no[k]:=l;
       end;
       {xac dinh doan trong danh sach moi ung voi tung chu cuoi}
       if t[i,length(t[i])]=b then inc(ee[b])
       else begin
          b:=t[i,length(t[i])]; bb[b]:=i; ee[b]:=i;
       end;
   end;
   close(f);
end;

procedure find(st,ed:longint; e:char);
{tim cac muc da biet vi tri ket thuc la ed, vi tri dau khong nho hon st,
 va ma tu chua trong do co chu cuoi e}
var i,k,l:integer;
    j:longint;
    can:boolean;
begin
   for k:=bb[e] to ee[e] do begin
       j:=ed; i:=length(t[k])-1;
       while (i>0) and (j>=st) do begin
          j:=last[j and max]^[t[k,i]]; dec(i);
       end;
       {lan luot tim cac chu cai cua ma tu}
       if j>=st then begin
          l:=m; can:=true;
          while (l>0) and (tt^[l]>=j) do begin
             if (ss^[l]>=j) and
                (length(t[nn^[l]])>=length(t[k])) then begin
                can:=false; break;
             end;
             dec(l);
          end;
          {chon loc cac muc; can=false nghia la muc moi kem muc da co}
          if can then begin
             inc(m); ss^[m]:=j; tt^[m]:=ed; nn^[m]:=k;
          end;
       end;
   end;
end;

procedure findanswer; {tim dap an toi uu}
var i,j,k:integer;
    f:text;
    len,next:^arr1;
    {len[i] la tong do dai ma tu trong dap an toi uu xet tren i muc dau
     (khong nhat thiet chon muc i)}
    {len[i]=max(len[j])+do dai ma tu cua muc i
     (i<j<=m, muc i va muc j khong co phan chung)}
    {next[i] ghi lua chon j trong cong thuc tren}
begin
   new(next); next^[0]:=0;
   new(len); len^[0]:=0;
   k:=0;
   for i:=1 to m do begin
      len^[i]:=len^[i-1]; {khong chon muc i}
      {tinh dap an toi uu neu chon muc i}
      j:=i-1;
      while (j>0) and (tt^[j]>=ss^[i]) do dec(j);
      if len^[j]+length(t[nn^[i]])>len^[i] then begin
         {xet xem chon muc i co tot hon khong}
         len^[i]:=len^[j]+length(t[nn^[i]]); next^[i]:=j;
      end;
   end;
   assign(f,st2); rewrite(f);
   writeln(f,len^[m]); {xuat tong gia tri cua dap an}
   k:=m; {tim cac muc co trong dap an}
   while k<>0 do begin
      while len^[k]=len^[k-1] do dec(k); {bang nhau thi khong chon muc k}
      writeln(f,no[nn^[k]],' ',ss^[k],' ',tt^[k]);
      k:=next^[k];
   end;
   close(f);
end;

begin
   readcod;
   new(ss); new(tt); new(nn);
   assign(f,st12);
   settextbuf(f,buf);
   reset(f);
   fillchar(nearest,sizeof(nearest),0);
   for i:=0 to max do begin
      new(last[i]); fillchar(last[i]^,sizeof(last[i]^),0);
   end;
   {khoi tao}
   while not seekeoln(f) do begin
      inc(now);
      last[now and max]^:=nearest;
      read(f,b); nearest[b]:=now; {cap nhat lan xuat hien cuoi cua b}
      if bb[b]<=ee[b] then {xet cac ma tu co chu cuoi b}
         if now>=1000 then find(now-999,now,b)
         else find(1,now,b);
         {xac dinh can duoi vi tri dau max(now-999,1), roi tim muc}
   end;
   close(f);
   findanswer;
end.
\end{verbatim}

\subsection{\texttt{boat.pas}: Bài đi thuyền, dùng cấu trúc cây nhị phân}

Định dạng vào của bài này: dòng đầu là tổng số người; mỗi dòng sau là họ và
tên của một người, cách nhau bởi một dấu cách. Họ dài không quá \(5\), tên
dài không quá \(10\). Tệp ra mỗi dòng cho biết những người ngồi trên một
thuyền; dòng cuối là tổng số thuyền cần dùng.

\begin{verbatim}
{$A+,B-,D+,E+,F-,G-,I-,L+,N-,O-,P-,Q-,R-,S-,T-,V+,X+}
{$M 65520,0,655360}
const Maxn=5000; {so nguoi toi da}
type Fnamearr=array[1..Maxn] of string[5];
     Gnamearr=array[1..Maxn] of string[10];
var Fname:^Fnamearr; {ghi ho cua moi nguoi}
    Gname:^Gnamearr; {ghi ten cua moi nguoi}
    visit:array[1..Maxn] of boolean; {visit[i] cho biet nguoi i da tham chua}
    total,n:integer; {total la so thuyen can dung, n la tong so nguoi}
    fo:text;

procedure init; {doc ten moi nguoi}
var f:text;
    st,t:string;
    i,j:integer;
begin
   new(Fname); new(Gname);
   write('Input file name:'); readln(st); assign(f,st); reset(f);
   readln(f,n);
   for i:=1 to n do begin
      readln(f,t);
      j:=pos(' ',t);
      Fname^[i]:=copy(t,1,j-1);
      delete(t,1,j);
      Gname^[i]:=t;
   end;
   close(f);
end;

procedure print(i,j:integer); {xuat viec i va j cung ngoi mot thuyen}
var t:integer;
begin
   t:=17-length(Fname^[i])-length(Gname^[i]);
   writeln(fo,Fname^[i],' ',Gname^[i],'':t,'-- ',Fname^[j],' ',Gname^[j]);
   inc(total);
end;

procedure try(father:integer; var root:integer; var rest:byte);
{dung cay con goc root va xuat phuong an di thuyen trong cay con}
{father cho biet root la con nao cua cha: 0 la con trai, 1 la con phai}
{rest cho biet sau khi xuat phuong an cho cay con co mot nguoi di rieng khong}
{rest=1 nghia la co, va nguoi di rieng co the doi len vi tri goc root}
var j,k:integer;
    Lrest,Rrest:byte;
    {Lrest, Rrest la gia tri rest cua cay con trai va phai cua root}
begin
   visit[root]:=true;
   j:=1;
   while (j<=n) and (visit[j] or (Fname^[root]<>Fname^[j])) do inc(j);
   {tim mot nguoi j cung ho voi root va chua tung tham}
   Lrest:=0;
   if j<=n then try(0,j,Lrest); {dung cay con goc j va xuat phuong an}
   k:=1;
   while (k<=n) and (visit[k] or (Gname^[root]<>Gname^[k])) do inc(k);
   {tim mot nguoi k cung ten voi root va chua tung tham}
   Rrest:=0;
   if k<=n then try(1,k,Rrest); {dung cay con goc k va xuat phuong an}
   {phan bo thuyen theo viec hai cay con con nguoi du hay khong}
   if (Lrest=1) xor (Rrest=1) then begin {chi mot cay con con nguoi du}
      if Lrest=1 then print(root,j) else print(root,k);
      {dua nguoi con du cung ngoi voi root}
      rest:=0;
   end
   else if (Lrest=1) and (Rrest=1) then begin {ca hai cay con deu con nguoi}
      {dua quan he giua root va cha de chon nguoi di cung root,
       roi doi nguoi con lai len vi tri goc root}
      if father=0 then begin
         print(root,k); root:=j;
      end
      else begin
         print(root,j); root:=k
      end;
      rest:=1;
   end
   else rest:=1;
end;

procedure findanswer; {tim moi cay trong rung, tuc moi thanh phan lien thong}
var i,j:integer;
    rest:byte;
begin
   assign(fo,'output.txt'); rewrite(fo);
   for j:=1 to n do
      if not visit[j] then begin {con nguoi chua tham thi o mot cay khac}
         i:=j;
         try(1,i,rest);
         {lay diem chua tham lam goc, dung cay va xuat phuong an trong cay}
         if rest=1 then begin {cuoi cung co con lai mot nguoi khong}
            writeln(fo,Fname^[i],' ',Gname^[i]);
            inc(total);
         end;
      end;
   writeln(fo,'The total of boats is ',total); {xuat tong so thuyen}
   close(fo);
end;

begin
   init;
   findanswer;
end.
\end{verbatim}

\subsection{\texttt{under.pas}: \textit{Underground City}, dùng cấu trúc lưu trữ liên kết}

Định dạng vào ra của chương trình theo chuẩn của đề thi.

\begin{verbatim}
{$A+,B-,D+,E+,F-,G-,I-,L+,N-,O-,P-,Q-,R-,S-,T-,V+,X+}
{$M 65520,0,655360}
uses undertpu;
const c:array[1..4,1..2] of shortint=((0,-1),(-1,0),(1,0),(0,1));
      {do doi toa do theo bon huong}
      d:array[1..4] of char=('S','W','E','N'); {chu cai ung voi bon huong}
      st1='under.inp'; {ten tep vao}
type aa=record
           x,y:shortint;
           next:integer;
        end;
     arr=array[0..10000] of aa;
var a:array[1..100,1..100] of char;
    {ban do thanh pho ngam; toa do (x,y) ung voi a[x,y].
     a[x,y]='O' nghia la khong co tuong, a[x,y]='W' nghia la co tuong}
    z:array[-100..100,-100..100] of char;
    {ban do moi gom thong tin da biet, (0,0) la vi tri xuat phat.
     'O','W' lan luot la khong co tuong va co tuong}
    {ngoai ra dung 'Q' va 'U' cho o chua biet; o 'U' co the tham do,
     con o 'Q' thi khong}
    uk,can:^arr;
    {dung mang de hien thuc danh sach lien ket; nut dau la phan tu chi so 0}
    {kieu nut: x,y la toa do, next la dia chi nut ke trong mang}
    {uk ghi cac vi tri hien co the tham do nhung chua biet trang thai;
     can ghi toa do cac diem xuat phat con kha nang}
    tail,rest,v,u:integer;
    {rest la so diem xuat phat con kha nang}
    {tail la con tro duoi cua danh sach uk, dong thoi cho biet uk da dung
     khong gian 0..tail-1}
    {u va v lan luot la so cot va so hang cua ban do}

procedure readp; {doc ban do thanh pho}
var f:text;
    i,j:integer;
begin
   assign(f,st1); reset(f);
   readln(f,u,v);
   for i:=v downto 1 do begin
      for j:=1 to u do read(f,a[j,i]);
      readln(f);
   end;
   close(f);
end;

procedure main; {thu tuc chinh tim diem xuat phat}
var nowx,nowy,nextx,nexty,o:integer;
    {vi tri hien tai so voi diem xuat phat la (nowx,nowy),
     muc tieu tham do tiep theo la (nextx,nexty)}
    e:char;

  procedure changeUK(xx,yy:integer);
  {sau khi them o khong co tuong (xx,yy),
   co the tham do cac o chua biet ke no}
  var i,j:integer;
  begin
     for i:=1 to 4 do
        if z[xx+c[i,1],yy+c[i,2]]='Q' then begin {xet bon o ke}
           z[xx+c[i,1],yy+c[i,2]]:='U'; {dat thanh co the tham do}
           with uk^[tail] do begin {them o co the tham do vao cuoi danh sach}
              x:=xx+c[i,1]; y:=yy+c[i,2];
              next:=tail+1;
           end;
           inc(tail); {de tail tro toi mot vi tri chua dung}
        end;
  end;

  procedure getready; {chuan bi truoc khi tham do}
  var i,j:integer;
  begin
     nowx:=0; nowy:=0;
     fillchar(z,sizeof(z),'Q');
     z[0,0]:='O'; {ban dau chi biet diem xuat phat khong co tuong}
     new(can); can^[0].next:=1;
     {thong ke moi o co the la diem xuat phat va dua vao danh sach can}
     for i:=1 to u do
        for j:=1 to v do
           if a[i,j]='O' then begin
              inc(rest);
              with can^[rest] do begin
                 x:=i; y:=j; next:=rest+1;
              end;
           end;
     can^[rest].next:=0;
     new(uk); tail:=1; uk^[0].next:=1;
     changeUK(0,0); {cac o ke diem xuat phat thanh o co the tham do}
  end;

  procedure change; {tim o xac dinh duoc ma khong can tham do}
  var px,py,i,j,d:integer;
      {d ghi trang thai: 1 la co tuong, 2 la khong co tuong,
       3 la ca hai deu co the}
  begin
     i:=0;
     while uk^[i].next<>tail do begin {duyet o chua biet co the tham do}
        px:=uk^[uk^[i].next].x; py:=uk^[uk^[i].next].y;
        d:=0; j:=0;
        while (d<3) and (can^[j].next<>0) do begin
           j:=can^[j].next;
           if a[can^[j].x+px,can^[j].y+py]='W'
              then d:=d or 1 {ghi kha nang co tuong}
              else d:=d or 2; {ghi kha nang khong co tuong}
        end;
        if d<3 then begin {neu khong phai ca hai trang thai deu co the}
           uk^[i].next:=uk^[uk^[i].next].next;
           {trang thai o nay da xac dinh, xoa khoi danh sach}
           if d=1 then z[px,py]:='W'; {chi co the la co tuong}
           if d=2 then begin {chi co the la khong co tuong}
              z[px,py]:='O';
              changeUK(px,py); {them cac o ke (px,py) co the tham do}
           end;
        end
        else i:=uk^[i].next;
     end;
  end;

  procedure find(var xx,yy:integer); {tim vi tri tham do tot nhat}
  var min,px,py,d1,d2,i,j,o:integer;
      {min la muc phu thuoc may man; cang nho cang tot}
  begin
     min:=maxint; i:=0;
     while uk^[i].next<>tail do begin
        px:=uk^[uk^[i].next].x; py:=uk^[uk^[i].next].y;
        d1:=0; j:=0;
        while can^[j].next<>0 do begin
           j:=can^[j].next;
           if a[can^[j].x+px,can^[j].y+py]='W'
              then inc(d1); {so kha nang co tuong tang 1}
        end;
        d2:=rest-d1; {tinh so kha nang khong co tuong}
        if abs(d1-d2)<min then begin {chon (px,py) neu can bang hon}
           min:=abs(d1-d2); o:=i; xx:=px; yy:=py;
        end;
        i:=uk^[i].next;
     end;
     uk^[o].next:=uk^[uk^[o].next].next; {xoa muc tieu tham do khoi danh sach}
  end;

  procedure path(x1,y1,x2,y2:integer);
  {tim phuong an di chuyen va tham do, roi loc diem xuat phat}
  var s,t,k,l,i,j:integer;
      p:arr;
  begin
     {dung BFS tim duong tu (x2,y2) den (x1,y1)}
     p[1].x:=x2; p[1].y:=y2; p[1].next:=0; s:=1; t:=1;
     while (p[s].x<>x1) or (p[s].y<>y1) do begin
        for k:=1 to 4 do
           if z[p[s].x+c[k,1],p[s].y+c[k,2]]='O' then begin
              inc(t);
              p[t].x:=p[s].x+c[k,1]; p[t].y:=p[s].y+c[k,2];
              z[p[t].x,p[t].y]:='P'; {danh dau vi tri da tim}
              p[t].next:=s; {ghi cha cua nut t}
           end;
        inc(s);
     end;
     {di chuyen den o ke vi tri can tham do (x2,y2)}
     l:=s; j:=p[l].next;
     while j<>1 do begin
        for i:=1 to 4 do
           if (p[l].x+c[i,1]=p[j].x) and
              (p[l].y+c[i,2]=p[j].y) then break;
        move(d[i]);
        l:=j; j:=p[l].next;
     end;
     {tim huong tham do}
     nowx:=p[l].x; nowy:=p[l].y;
     for i:=1 to 4 do
        if (p[l].x+c[i,1]=p[1].x) and
           (p[l].y+c[i,2]=p[1].y) then break;
     e:=look(d[i]); {tham do}
     z[x2,y2]:=e; {danh dau tren ban do}
     if e='O' then changeUK(x2,y2); {neu khong co tuong, them vi tri}
     for i:=-u to u do
        for j:=-v to v do
           if z[i,j]='P' then z[i,j]:='O'; {khoi phuc ban do}
     j:=0;
     while can^[j].next<>0 do begin {loc diem xuat phat}
        k:=j;
        j:=can^[j].next;
        if a[can^[j].x+x2,can^[j].y+y2]<>e then begin
           {trang thai cua vi tri tuong doi khong khop ket qua tham do}
           dec(rest); can^[k].next:=can^[j].next; j:=k;
        end;
     end;
  end;

begin
   getready; {chuan bi truoc khi tham do}
   start; {bat dau tham do}
   while rest>1 do begin
      change; {tim o co trang thai xac dinh ma khong can tham do}
      find(nextx,nexty); {tim vi tri tham do tot nhat}
      path(nowx,nowy,nextx,nexty);
      {tim qua trinh di chuyen, tham do cu the va loc diem xuat phat}
   end;
   o:=can^[0].next;
   finish(can^[o].x,can^[o].y); {tim duoc vi tri xuat phat}
end;

begin
   readp;
   main;
end.
\end{verbatim}

\subsection{\texttt{sequence.pas}: Dãy nhỏ nhất, kết hợp cấu trúc tuyến tính và đống nhị phân}

Dữ liệu vào của chương trình: dòng đầu là cạnh \(N\) của ma trận; tiếp theo là
\(N\) dòng chứa ma trận số nguyên dương \(N\times N\); dòng cuối gồm bốn số,
biểu diễn vị trí đầu và vị trí cuối theo thứ tự hàng rồi cột. Tệp ra có dòng
đầu là giá trị nhỏ nhất của tổng trị tuyệt đối hiệu giữa hai số kề nhau trong
dãy, dòng thứ hai là dãy kề nhau tương ứng.

\begin{verbatim}
{$A+,B-,D+,E+,F-,G-,I+,L+,N-,O-,P-,Q-,R-,S+,T-,V+,X+}
{$M 16384,0,655360}
const b:array[1..4,1..2] of shortint=((1,0),(0,1),(-1,0),(0,-1));
type arr=array[0..101,0..101] of integer;
     aa=record
        {kieu nut trong dong; w la khoang cach den dich,
         (x,y) la vi tri cua nut trong ma tran}
        w:integer;
        x,y:byte;
     end;
var a,d:array[0..101,0..101] of byte;
    {a[i,j] ghi ma tran; d[i,j] la huong buoc tiep theo trong duong ngan nhat
     tu (i,j) den dich}
    pos:^arr;
    {pos[i,j] la mang anh xa giua cac cau truc du lieu, bieu thi vi tri luu
     khoang cach tu (i,j) den dich trong cau truc dong, tuc h[pos[i,j]].w}
    h:array[0..10001] of aa;
    {dong nhi phan bieu dien bang mang; con trai cua h[i] la h[i*2],
     con phai la h[i*2+1]}
    x1,y1,x2,y2,n:byte; {n la canh ma tran; (x1,y1),(x2,y2) la dau va dich}
    t:integer; {tong so nut trong dong}

procedure readp; {doc du lieu}
var f:text;
    st:string;
    i,j:byte;
begin
   new(pos);
   write('File name:'); readln(st); assign(f,st); reset(f);
   readln(f,n);
   for i:=1 to n do
      for j:=1 to n do
         read(f,a[i,j]);
   readln(f,x1,y1,x2,y2);
   close(f);
end;

procedure swap(i,j:integer); {doi nut thu i va thu j trong dong}
var k:aa;
begin
   k:=h[i]; h[i]:=h[j]; h[j]:=k; {doi khoang cach}
   pos^[h[i].x,h[i].y]:=i; pos^[h[j].x,h[j].y]:=j; {doi vi tri anh xa}
end;

procedure heap1(i:integer); {vun dong tu nut i xuong duoi}
var j:integer;
begin
   while i*2<=t do begin
      if (i*2+1<=t) and (h[i*2+1].w<h[i*2].w)
         {tim con co khoang cach den dich ngan hon}
         then j:=i*2+1
         else j:=i*2;
      if h[j].w<h[i].w then begin {co can doi voi nut goc khong}
         swap(i,j);
         i:=j;
      end
      else break;
   end;
end;

procedure heap2(i:integer); {vun dong tu nut i len tren}
begin
   while (i<>1) and (h[i div 2].w>h[i].w) do begin {co can doi voi cha khong}
      swap(i div 2,i);
      i:=i div 2;
   end;
end;

procedure main; {tim duong ngan nhat tu moi vi tri den dich}
var xx,yy,i,j:byte;
    fo:text;
begin
   h[10001].w:=9999;
   for i:=1 to n do
      for j:=1 to n do
         pos^[i,j]:=10001;
   {dat khoang cach tu moi vi tri den dich la vo cuc}
   h[0].w:=0;
   for i:=1 to n do begin
      pos^[i,0]:=0; pos^[0,i]:=0; pos^[i,n+1]:=0; pos^[n+1,i]:=0;
   end;
   {them mot vong quanh ma tran de khong di ra ngoai, gia su khoang cach
    den vong nay deu la 0}
   {vi 0 khong the nho hon nua, se khong di ra ngoai ma tran}
   with h[1] do begin
      w:=0; x:=x2; y:=y2;
   end;
   pos^[x2,y2]:=1; t:=1; {dua dich lam goc dong}
   repeat
      {nut goc trong dong la diem dau cua duong ngan nhat}
      for i:=1 to 4 do begin {doi khoang cach tu cac diem ke goc den dich}
         xx:=h[1].x+b[i,1]; yy:=h[1].y+b[i,2]; {tinh toa do diem ke}
         if h[pos^[xx,yy]].w=9999 then begin {xet da o trong dong chua}
            inc(t);
            with h[t] do begin {them mot nut moi vao dong}
               w:=h[1].w+abs(a[h[1].x,h[1].y]-a[xx,yy]);
               x:=xx; y:=yy;
               pos^[xx,yy]:=t; d[xx,yy]:=i;
            end;
            heap2(t); {vun len tu nut moi them}
         end
         else if h[1].w+abs(a[h[1].x,h[1].y]-a[xx,yy])<h[pos^[xx,yy]].w
              then begin {xet co tim duoc khoang cach ngan hon khong}
                 h[pos^[xx,yy]].w:=h[1].w+abs(a[h[1].x,h[1].y]-a[xx,yy]);
                 {cap nhat khoang cach}
                 heap2(pos^[xx,yy]); {vun len tu diem tuong ung trong dong}
                 d[xx,yy]:=i;
              end;
      end;
      swap(1,t); pos^[h[t].x,h[t].y]:=0; dec(t); {lay nut goc ra}
      heap1(1); {vun xuong tu nut goc}
   until (h[1].x=x1) and (h[1].y=y1); {da tim duong tu dau den dich}
   assign(fo,'output.txt'); rewrite(fo);
   writeln(fo,h[1].w); {xuat khoang cach}
   {xuat cac diem tren duong di}
   while (x1<>x2) or (y1<>y2) do begin
      write(fo,a[x1,y1],' '); i:=x1; j:=y1;
      dec(x1,b[d[i,j],1]); dec(y1,b[d[i,j],2]);
   end;
   writeln(fo,a[x2,y2]);
   close(fo);
end;

begin
   readp;
   main;
end.
\end{verbatim}

\section{Tài liệu tham khảo}

\begin{enumerate}
  \item Yan Weimin, Wu Weimin. \textit{Cấu trúc dữ liệu}, bản thứ hai. Nhà
  xuất bản Đại học Thanh Hoa.
  \item Wu Wenhu, Wang Jiande. \textit{Phân tích và thiết kế chương trình cho
  các thuật toán thực dụng}. Nhà xuất bản Công nghiệp Điện tử.
  \item Wu Wenhu, Wang Jiande. \textit{Hướng dẫn thi Olympic Tin học quốc tế
  và toàn quốc cho thanh thiếu niên: thuật toán đồ thị và lập trình}.
  \item \textit{Olympic Tin học}, tạp chí quý, số 1 và số 2 năm 1998.
  \item Đề IOI 1999 và các đề thi cấp tỉnh Hồ Nam qua các năm.
\end{enumerate}

\end{document}
